\section{Alternative Adaptation Strategies}
\label{sec:baselines}
% Promoted from Background 2.4 on 2026-09-17. These are the arms the merge is compared against; the
% six-arm figure sits with them.

\begin{figure*}[t]
\centering
\begin{tikzpicture}[
  font=\scriptsize,
  base/.style ={draw, rounded corners=2pt, minimum height=4.6mm, minimum width=12mm, fill=black!5},
  ad/.style   ={draw, rounded corners=1pt, minimum height=3.6mm, minimum width=7mm, fill=blue!10},
  adw/.style  ={ad, fill=green!18},
  gt/.style   ={draw, diamond, aspect=2.4, inner sep=0.8pt, fill=orange!20, minimum height=3.6mm},
  clab/.style  ={font=\scriptsize, align=center, text depth=0pt, text width=34mm},
  ar/.style   ={-{Latex[length=1.2mm]}, shorten >=0.5pt},
]
\def\cx{52mm}\def\ry{28mm}
% ---------- row 1: no adaptation, then adaptation without composition --------------------------
\begin{scope}[shift={(0,0)}]
  \node[base] (b1) {base};
  \node[clab, below=1.4mm of b1] {\textbf{untuned}\\[-1pt]no adaptation};
\end{scope}
\begin{scope}[shift={(\cx,0)}]
  \node[base] (b2) {base};
  \node[draw, dashed, rounded corners=1pt, minimum height=3.6mm, minimum width=10mm,
        above=1.8mm of b2, fill=yellow!22] (ex) {example};
  \draw[ar] (ex) -- (b2);
  \node[clab, below=1.4mm of b2] {\textbf{in-context}\\[-1pt]one solved example};
\end{scope}
\begin{scope}[shift={(2*\cx,0)}]
  \node[base] (b3) {base};
  \node[ad, above=1.8mm of b3] (a3) {\texttt{L0}};
  \draw[ar] (a3) -- (b3);
  \node[clab, below=1.4mm of b3] {\textbf{clean LoRA}\\[-1pt]clean code only};
\end{scope}
% ---------- row 2: the three ways to use the per-transform data --------------------------------
% All three bottom panels now show the SAME per-condition objects, labelled by obfuscator
% (2026-09-20, user): pooled's inputs were unlabelled grey squares while the router's and the
% merge's were blue adapter boxes, so the figure read as three different inputs when the paper's
% whole argument is that they are the same six conditions combined three ways. Pooled's squares are
% its training DATA rather than adapters, which is why they keep a distinct shape -- flat and wider
% -- but they now carry the same fill and the same labels.
\begin{scope}[shift={(0,-\ry)}]
  \node[base] (b4) {base};
  \node[ad, above=1.8mm of b4, minimum width=13mm, fill=red!14] (a4) {pooled};
  \draw[ar] (a4) -- (b4);
  \foreach \i/\x/\t in {1/-9mm/{\tiny L1b}, 2/0mm/{\tiny S1}, 3/9mm/{\tiny$\cdots$}}
    \node[ad, minimum width=7mm, minimum height=3mm, above=2.6mm of a4, xshift=\x] (p\i) {\t};
  \foreach \i in {1,2,3} \draw[ar] (p\i) -- (a4);
  \node[clab, below=1.4mm of b4] {\textbf{pooled}\\[-1pt]rows merged, one adapter};
\end{scope}
\begin{scope}[shift={(\cx,-\ry)}]
  \node[base] (b5) {base};
  \node[gt, above=1.8mm of b5] (g5) {gate};
  \draw[ar] (g5) -- (b5);
  \foreach \i/\x/\t in {1/-9mm/{\tiny L1b}, 2/0mm/{\tiny S1}, 3/9mm/{\tiny$\cdots$}}
    \node[ad, above=2mm of g5, xshift=\x, minimum width=7mm, minimum height=3.6mm] (e\i) {\t};
  % ROUTING SHOWN, NOT JUST MIXING (2026-09-20, user): the gate selects. Solid paths are the experts
  % this input routes to; the dashed path is an expert the gate leaves out. Without this the panel
  % looked identical to the merge's, which combines all of them unconditionally.
  \draw[ar] (e1) -- (g5);
  \draw[ar] (e2) -- (g5);
  \draw[ar, dashed, black!55] (e3) -- (g5);
  \node[font=\tiny, text=black!55, anchor=west] at ($(e3)!0.5!(g5) + (1.2mm,0.6mm)$) {not routed};
  \node[clab, below=1.4mm of b5] {\textbf{router}\\[-1pt]gate routes per input};
\end{scope}
\begin{scope}[shift={(2*\cx,-\ry)}]
  \node[base] (b6) {base};
  \node[adw, above=1.8mm of b6, minimum width=13mm] (m6) {merged};
  \draw[ar] (m6) -- (b6);
  \foreach \i/\x/\t in {1/-9mm/{\tiny L1b}, 2/0mm/{\tiny S1}, 3/9mm/{\tiny$\cdots$}}
    \node[ad, above=2.6mm of m6, xshift=\x, minimum width=7mm, minimum height=3.6mm] (f\i) {\t};
  \foreach \i in {1,2,3} \draw[ar] (f\i) -- (m6);
  \node[clab, below=1.4mm of b6] {\textbf{merge}\\[-1pt]combined beforehand};
\end{scope}
% The three grey axis captions were removed 2026-09-20: the caption already states the ordering,
% and the labels repeated it in two words each.
\end{tikzpicture}
\caption{\textbf{The six arms.} All are the same base model at the same LoRA rank, differing only in
how the training data is used and where per-transform knowledge is combined. Top row: arms that never
see an obfuscated program in training --- \emph{untuned}, \emph{in-context}, and the \emph{clean
LoRA} control that separates learning the task from learning the transforms. Bottom row: the three ways to spend the obfuscated data. The labelled boxes are the six training
conditions, named by obfuscator (\texttt{L1b}, \texttt{S1}, and four more); all three arms receive
the same six. In the \emph{router} panel the dashed path marks an expert the gate does not select for a given input; the merge combines all of them unconditionally. The three are ordered by \emph{when} the conditions are combined ---
\emph{pooled} collapses them during training, \emph{router} composes them per input at inference and
must keep every expert resident, and \emph{merge} composes them in weight space beforehand, leaving
one adapter with no gate and no run-time cost. The KL-consistency arm is \emph{pooled} with an added
consistency term (Section~\ref{sec:baselines}).}
\label{fig:baselines}
\end{figure*}

Figure~\ref{fig:baselines} places the merge among the five other ways one might spend the same budget
of obfuscated training data. The first two do not adapt weights at all and bound the problem; the next
three are the standard single-adapter answers; the last composes. Every arm is a LoRA of the same rank
on the same target modules of the same base model, so what differs is the data and the combination
rule, not the capacity. We give each enough detail to be reimplemented; exact hyperparameters are in
Section~\ref{sec:setup}.

\noindent\textbf{Untuned.} The instruction-tuned checkpoint, prompted directly, with no weight change.
It is the reference every other arm is measured against, and on the backward task it is a surprisingly
strong one.

\noindent\textbf{In-context.} The same checkpoint with one solved example prepended to the
prompt~\cite{fewshot}. The example is a single frozen seven-line program written for the purpose ---
not drawn from the corpus --- shown \emph{unobfuscated} whatever condition is being evaluated, with
its gold value obtained by executing it through the same sandbox as every other item. It is identical
for every item, model and condition, and its identity is recorded in the prompt-template hash, so no
result here depends on example selection. Because the example is clean code, this arm sees no
obfuscated program at any point: it isolates whether the \emph{answer format}, rather than the
obfuscation, is what the untuned model is failing at, and it is the weakest intervention in the study
by design.

\noindent\textbf{Clean-code LoRA.} An adapter trained on unobfuscated programs only, with the same
prompt, optimiser and step count as every other adapter. It is the control that separates
\emph{learning the task} from \emph{learning the transforms}: it sees the prompt format and the answer
distribution but never an obfuscated program, so whatever it gains is not transform knowledge. It is
also the teacher for the KL arm below.

\noindent\textbf{Pooled training.} One adapter trained on the union of every training condition,
shuffled. This is the obvious way to spend the data and the one most likely to be done in practice: it
needs no per-transform bookkeeping, no gate and no merge step. Because every composed arm below is
built from specialists trained on exactly these rows, pooled training is the data-matched control for
all of them.

\noindent\textbf{KL-consistency training.} Pooled training plus a distillation
term~\cite{distillation} that ties the student's answer to its own behaviour on the clean parent. For
an obfuscated program $x_{\mathrm{obf}}$ with clean parent $x_{\mathrm{clean}}$ and gold answer $y$,
%
\begin{equation}
\mathcal{L} \;=\; \mathrm{CE}\!\left(y \mid x_{\mathrm{obf}}\right)
\;+\; \lambda\, D_{\mathrm{KL}}\!\left(
p_{\mathrm{teacher}}(\cdot \mid x_{\mathrm{clean}}) \,\Vert\,
p_{\mathrm{student}}(\cdot \mid x_{\mathrm{obf}}) \right),
\label{eq:kl}
\end{equation}
%
where the teacher is the frozen clean-code LoRA above, the KL is taken over the answer tokens only,
and the student never sees the parent. Three properties make it a fair comparison rather than an extra
knob. It trains on exactly the rows pooled training uses, for the same number of steps. At
$\lambda = 0$ it reduces to pooled training identically, so pooled training is not an approximate
control but an exact one. And the teacher is fixed across all conditions, so the term adds no
transform-specific supervision --- it only asks that the same computation, written two ways, be
answered the same way.

\noindent\textbf{Router system.} One adapter per transform, kept separate, with a learned gate that
mixes their low-rank updates per token and per layer. The gate reads the decoder layer's input hidden state $h$, so one gate serves every target
projection in that layer. It scores experts by attention against a learned key $k_e$ per expert:
%
\begin{equation}
a \;=\; \mathrm{softmax}_e\!\left( \frac{ \left(W_q\,\mathrm{LayerNorm}(h)\right) \cdot k_e }{ \sqrt{d_r}\,\tau } \right),
\label{eq:gate}
\end{equation}
%
with router width $d_r$ and a learned temperature $\tau$ per layer. Only the gate is trained; the
experts and the base model stay frozen, so the arm adds a small number of parameters relative to the
expert bank. A learned gate may be worth nothing over no gate. We therefore report three controls: the same gate
argmaxed to a hard route, a fixed uniform gate, and a gate frozen at random initialisation. These
separate the value of the mixture from the value of the routing. This arm is the most expressive of the six and the most expensive: every
expert must be resident at inference, and the gate runs at every layer for every token.
